\section{System Architecture and Design}
\label{sec:architecture}

The Mini-A-IPC framework is architected as a modular, three-tier systems evaluation and adaptive execution platform. As illustrated in Fig.~\ref{fig:architecture}, the system decouples high-level workload orchestration from low-level communication primitives while providing seamless integration between user-space mechanisms and custom kernel modules.

\begin{figure}[t]
\centering
\begin{tikzpicture}[
    scale=0.82, every node/.style={transform shape},
    layer/.style={draw=neutralgrey, dashed, fill=neutralgrey!5, inner sep=6pt, rounded corners=3pt},
    banner/.style={fill=darknavy, text=white, font=\scriptsize\bfseries, rounded corners=2pt, inner sep=2.5pt},
    arr/.style={->, >=stealth, thick, draw=charcoal},
    ctrlarr/.style={->, >=stealth, thick, draw=acadblue},
    dasharr/.style={->, >=stealth, thick, dashed, draw=mutedwine}
]

% =========================================================================
% REUSABLE GRID, DIMENSION, AND SPACING CONSTANTS
% =========================================================================
\def\boxW{2.40}         % Exact uniform width for ALL 13 component boxes in Areas A, B, C (cm)
\def\boxH{0.95}         % Exact uniform height for ALL 13 component boxes in Areas A, B, C (cm)
\def\rightW{5.30}       % Exact uniform width for Right-Side Adaptive & Benchmarking boxes (cm)

\def\gridPitchX{2.80}   % Strict horizontal grid spacing (center-to-center pitch)
\def\xOne{1.40}         % Column 1 Center X
\def\xTwo{4.20}         % Column 2 Center X (\xOne + \gridPitchX)
\def\xThree{7.00}       % Column 3 Center X (\xTwo + \gridPitchX)
\def\xFour{9.80}        % Column 4 Center X (\xThree + \gridPitchX)
\def\xMidTwoThree{5.60} % Midpoint between Column 2 and Column 3 (whitespace corridor)
\def\xCorridor{8.40}    % Midpoint between Column 3 right edge (8.20) and Col 4 left edge (8.60)
\def\xRight{15.65}      % Center X for Adaptive & Benchmarking suite (Left edge at 13.00, Corridor = 2.00cm!)

\def\yApp{8.45}         % Area A: Application & API Layer Y-center
\def\yBus{7.05}         % Workload Dispatch Bus Rail Y-coordinate (Spacious separation from Banner B!)
\def\yTransOne{5.95}    % Area B: Transport Layer Row 1 Y-center
\def\yTransTwo{4.35}    % Area B: Transport Layer Row 2 Y-center
\def\yBound{2.85}       % User-Kernel Execution Boundary Y-coordinate
\def\yKern{1.40}        % Area C: Kernel Execution Layer Y-center
\def\yAdapt{6.25}       % Area D: Adaptive Transport Selector Y-center
\def\yBench{1.70}       % Area D: Benchmarking & Telemetry Suite Y-center

% Standardized Component Box Styles (Strict Uniformity Across All Layers)
\tikzset{
    box/.style={
        draw=darknavy, rounded corners=2.5pt, fill=medblue!12, thick,
        minimum width=\boxW cm, text width=\boxW cm - 6pt,
        minimum height=\boxH cm,
        align=center, font=\footnotesize, inner sep=2.5pt
    },
    apibox/.style={
        draw=darknavy, rounded corners=2.5pt, fill=acadblue!15, thick,
        minimum width=\boxW cm, text width=\boxW cm - 6pt,
        minimum height=\boxH cm,
        align=center, font=\footnotesize, inner sep=2.5pt
    },
    kernbox/.style={
        draw=darknavy, rounded corners=2.5pt, fill=mutedorange!15, thick,
        minimum width=\boxW cm, text width=\boxW cm - 6pt,
        minimum height=\boxH cm,
        align=center, font=\footnotesize, inner sep=2.5pt
    },
    kmodbox/.style={
        draw=darknavy, rounded corners=2.5pt, fill=mutedorange!25, thick,
        minimum width=\boxW cm, text width=\boxW cm - 6pt,
        minimum height=\boxH cm,
        align=center, font=\footnotesize, inner sep=2.5pt
    },
    adaptbox/.style={
        draw=darknavy, rounded corners=2.5pt, fill=mutedgreen!15, thick,
        minimum width=\rightW cm, text width=\rightW cm - 12pt,
        minimum height=3.05cm,
        align=center, font=\footnotesize, inner sep=5pt
    },
    benchbox/.style={
        draw=darknavy, rounded corners=2.5pt, fill=acadblue!12, thick,
        minimum width=\rightW cm, text width=\rightW cm - 12pt,
        minimum height=2.40cm,
        align=center, font=\footnotesize, inner sep=5pt
    }
}

% =========================================================================
% AREA A: APPLICATION AND API LAYER (Identical Height, Aligned Tops & Widths)
% =========================================================================
\node[apibox] (cli) at (\xOne, \yApp) {
    \textbf{CLI Applications}\\[2pt]
    {\scriptsize \texttt{bin/mini-aipc} $\cdot$ \texttt{client}}
};

\node[apibox] (iface) at (\xTwo, \yApp) {
    \textbf{Unified IPC API}\\[2pt]
    {\scriptsize \texttt{include/ipc.h} (\texttt{ipc\_opts})}
};

\node[apibox] (disp) at (\xThree, \yApp) {
    \textbf{Workload Dispatcher}\\[2pt]
    {\scriptsize Single RTT $\cdot$ MPMC $\cdot$ Stream}
};

% Application Layer Horizontal Workload Flow
\draw[arr] (cli) -- (iface);
\draw[arr] (iface) -- (disp);

% Feeder from Workload Dispatcher into Dispatch Bus (Enters bus centrally at x = \xMidTwoThree)
\draw[thick, draw=charcoal] (disp.south) -- ++(0,-0.60) -| (\xMidTwoThree, \yBus);

% =========================================================================
% AREA B: IPC TRANSPORT LAYER (Strict 3-Column x 2-Row Grid, Uniform Boxes)
% =========================================================================

% Workload Dispatch Bus Rail (Above Row 1 at y = \yBus = 7.05)
\draw[thick, draw=charcoal] (\xOne, \yBus) -- (\xThree, \yBus);
\node[above=2pt, font=\scriptsize\bfseries, text=darknavy, fill=white, inner sep=1.5pt, rounded corners=2pt] at (3.80, \yBus) {Workload Dispatch Bus};

% Column 1: Stream Transports (x = \xOne)
\node[box] (pipe) at (\xOne, \yTransOne) {
    \textbf{Anonymous Pipes}\\[3pt]
    {\tiny \texttt{src/ipc\_pipe.c} $\cdot$ \texttt{pipe2()}}
};
\node[box] (fifo) at (\xOne, \yTransTwo) {
    \textbf{Named FIFOs}\\[3pt]
    {\tiny \texttt{src/ipc\_fifo.c} $\cdot$ \texttt{mkfifo()}}
};

% Column 2: Message Queues (x = \xTwo)
\node[box] (sysv) at (\xTwo, \yTransOne) {
    \textbf{System V MQ}\\[3pt]
    {\tiny \texttt{src/ipc\_sysv\_mq.c} $\cdot$ \texttt{msgget()}}
};
\node[box] (posix) at (\xTwo, \yTransTwo) {
    \textbf{POSIX MQ}\\[3pt]
    {\tiny \texttt{src/ipc\_posix\_mq.c} $\cdot$ \texttt{mq\_open()}}
};

% Column 3: Direct Memory & Sockets (x = \xThree)
\node[box] (shm) at (\xThree, \yTransOne) {
    \textbf{POSIX Shared Memory}\\[3pt]
    {\tiny \texttt{src/ipc\_shm.c} $\cdot$ Futex Mailbox}
};
\node[box] (uds) at (\xThree, \yTransTwo) {
    \textbf{Unix Domain Sockets}\\[3pt]
    {\tiny \texttt{src/ipc\_uds.c} $\cdot$ \texttt{AF\_UNIX}}
};

% Dedicated External Branch Connectors from Bus to Top Edges of Row 1 Transports
\draw[arr] (\xOne, \yBus) -- (pipe.north);
\draw[arr] (\xTwo, \yBus) -- (sysv.north);
\draw[arr] (\xThree, \yBus) -- (shm.north);

% Column drops within stream and queue families
\draw[arr] (pipe.south) -- (fifo.north);
\draw[arr] (sysv.south) -- (posix.north);

% Dedicated External Dispatch Branch to Unix Domain Sockets (Enters UDS at west port, zero SHM crossing)
\draw[arr] (\xMidTwoThree, \yBus) -- (\xMidTwoThree, \yTransTwo) -- (uds.west);

% =========================================================================
% USER-KERNEL EXECUTION BOUNDARY (Dedicated Whitespace, Zero Overlap)
% =========================================================================
\draw[thick, dashed, draw=neutralgrey] (-0.20, \yBound) -- (11.20, \yBound);
\node[font=\tiny\bfseries, fill=white, draw=neutralgrey, rounded corners=2pt, inner sep=2pt, text=darknavy] at (\xMidTwoThree, \yBound) {USER--KERNEL BOUNDARY};

% =========================================================================
% AREA C: KERNEL EXECUTION LAYER (Strict 4-Box Horizontal Grid, Uniform Dimensions)
% =========================================================================
\node[kernbox] (vfs) at (\xOne, \yKern) {
    \textbf{Linux VFS Core}\\[2pt]
    {\scriptsize Pipe Page Buffers $\cdot$ \texttt{struct file}}
};

\node[kernbox] (mq) at (\xTwo, \yKern) {
    \textbf{Kernel MQs}\\[2pt]
    {\scriptsize SysV \texttt{ipc/msg.c} $\cdot$ \texttt{/dev/mqueue}}
};

\node[kernbox] (futex) at (\xThree, \yKern) {
    \textbf{Futex / MM / Skt}\\[2pt]
    {\scriptsize Zero-Copy SHM $\cdot$ \texttt{sk\_buff}}
};

\node[kmodbox] (kmod) at (\xFour, \yKern) {
    \textbf{Custom Mailbox}\\[2pt]
    {\scriptsize \texttt{/dev/aipc} $\cdot$ 64KB Ring Buffer}
};

% =========================================================================
% DEDICATED KERNEL CONNECTORS (Strictly Orthogonal, Zero Component Crossing)
% =========================================================================

% Column 1 -> Linux VFS Core: Direct vertical connection from Named FIFOs
\draw[arr] (fifo.south) -- (vfs.north);

% Column 2 -> Kernel MQs: Direct vertical connection from POSIX MQ
\draw[arr] (posix.south) -- (mq.north);

% Column 3 -> Futex / MM & Socket Core:
% Unix Domain Sockets connects straight down at center with masked label
\draw[arr] (uds.south) -- node[right=2pt, font=\tiny, text=charcoal, fill=white, inner sep=1pt] {AF\_UNIX} (futex.north);

% POSIX Shared Memory routes via external orthogonal bypass corridor around UDS
\draw[arr] ([yshift=-2mm]shm.east) -- (\xCorridor, 5.75) -- (\xCorridor, 2.25) node[midway, right=2pt, font=\tiny\bfseries, text=darknavy, fill=white, inner sep=1pt] {Futex/MM Bypass} -- (7.50, 2.25) -- ([xshift=5mm]futex.north);

% Column 4 -> Custom Mailbox /dev/aipc: Direct character device entry below boundary
\draw[arr] (\xFour, \yBound) -- node[right=2pt, font=\tiny, text=charcoal, fill=white, inner sep=1pt] {\texttt{/dev/aipc}} (kmod.north);

% =========================================================================
% AREA D: ADAPTIVE SELECTION & TELEMETRY SUITE (Uniform Width & Outer Geometry)
% =========================================================================

\node[adaptbox] (adapt) at (\xRight, \yAdapt) {
    \textbf{\normalsize Adaptive Transport Selector}\\[1.5pt]
    \scriptsize \texttt{src/selector.c}\\[4pt]
    \raggedright
    \textbf{1) Static Decision Policy (Thresholding):}\\[1.5pt]
    \quad $\bullet$ $s \le 1024\text{ B} \implies$ Route to \textbf{POSIX SHM}\\[1.5pt]
    \quad $\bullet$ $s > 1024\text{ B} \implies$ Route to \textbf{Anonymous Pipe}\\[4pt]
    \textbf{2) Dynamic EWMA Telemetry Controller:}\\[1.5pt]
    \quad $\bullet$ $\bar{L}^{(t)} = \alpha L + (1-\alpha)\bar{L}^{(t-1)}$ ($\alpha=0.15$)\\[1.5pt]
    \quad $\bullet$ 8\% Hysteresis Switching Margin ($h=1.08$)\\[3.5pt]
    \centering \textit{\scriptsize [Evaluated Candidates: SHM $\longleftrightarrow$ Pipe]}
};

\node[benchbox] (bench) at (\xRight, \yBench) {
    \textbf{\normalsize Benchmarking \& Telemetry Suite}\\[1.5pt]
    \scriptsize \texttt{bench/unified\_bench.c} $\cdot$ \texttt{orchestrate.py}\\[4pt]
    \raggedright
    $\bullet$ \textbf{Scientific Controls:} Fisher-Yates trial shuffle, dedicated CPU affinity, 500 untimed priming cycles\\[2.5pt]
    $\bullet$ \textbf{Telemetry Engine:} \texttt{CLOCK\_MONOTONIC} (vDSO ns), \texttt{getrusage} context switches, \texttt{strace} profiling
};

% =========================================================================
% INTER-MODULE CONTROL & TELEMETRY LINKS (Non-Crossing Dedicated Routes)
% =========================================================================

% Workload Dispatcher activates Adaptive Mode
\draw[ctrlarr] (disp.east) -- node[above=2pt, font=\scriptsize\bfseries, text=acadblue, fill=white, inner sep=1.5pt] {Adaptive Mode} (adapt.155);

% Adaptive Selector routes to Transport Layer (Clean horizontal link to upper SHM port)
\coordinate (shmport) at ([yshift=2mm]shm.east);
\coordinate (adaptport) at (adapt.west |- shmport);
\draw[ctrlarr] (adaptport) -- node[above=2pt, font=\scriptsize\bfseries, text=acadblue, fill=white, inner sep=1.5pt] {Select Route (SHM/Pipe)} (shmport);

% Internal Telemetry Feedback within Area D (Positioned cleanly midway between boxes)
\draw[dasharr] (bench.north) -- (adapt.south);
\path (bench.north) -- (adapt.south) coordinate[midway] (feedmid);
\node[right=3pt, font=\scriptsize\bfseries, text=mutedwine] at (feedmid) {EWMA Latency Feedback};

% OS Telemetry collection from Kernel Modules (Clean Horizontal Link Across Generous 2.00cm Whitespace Corridor!)
\coordinate (kmodport) at (kmod.east);
\coordinate (benchport) at (bench.west |- kmodport);
\draw[dasharr] (kmodport) -- (benchport);
\node[above=2pt, font=\scriptsize\bfseries, text=mutedwine, fill=white, inner sep=1.5pt] at (12.00, 1.40) {Kernel Telemetry};

% =========================================================================
% SECTION BANNERS & FRAMES
% =========================================================================

\node[banner, anchor=south west] at (-0.15, 9.15) {A. Application and API Layer};
\node[banner, anchor=south west] at (-0.15, 7.55) {B. Classical IPC Transport Layer};
\node[banner, anchor=south west] at (-0.15, 0.15) {C. Kernel Execution Layer \& Standalone Driver Mailbox};
\node[banner, anchor=south west] at (13.00, 9.15) {D. Adaptive Selection \& Benchmarking Suite};

\begin{scope}[on background layer]
    \node[layer, fit=(cli) (disp), inner ysep=6pt, inner xsep=5pt] {};
    \node[layer, fit=(pipe) (uds), inner ysep=6pt, inner xsep=5pt] {};
    \node[layer, fit=(vfs) (kmod), inner ysep=6pt, inner xsep=5pt] {};
    \node[layer, fit=(adapt) (bench), inner ysep=6pt, inner xsep=5pt] {};
\end{scope}

\end{tikzpicture}

\caption{Overall Mini-A-IPC Architectural Topology, illustrating the relationship between the userspace classical transports, the adaptive telemetry controller, the custom kernel mailbox device, and the benchmarking framework.}
\label{fig:architecture}
\end{figure}

\subsection{Overall Architecture}
The Mini-A-IPC architecture comprises four primary subsystems:
(1)~\textbf{Userspace Communication Layer:} Self-contained, POSIX-compliant implementations of the six classical Linux IPC mechanisms adhering to the standardized interface in \texttt{include/ipc.h};
(2)~\textbf{Adaptive Decision Layer:} Static size-threshold routing logic and a closed-loop EWMA telemetry controller (\texttt{src/selector.c}) that dynamically selects optimal transports;
(3)~\textbf{Kernel-Space Device Subsystem:} A custom Linux character device driver \texttt{/dev/aipc} (\texttt{kernel/aipc\_module.c}) providing in-kernel ring buffering with sleepable synchronization; and
(4)~\textbf{Benchmarking and Profiling Suite:} An automated C and Python benchmarking harness executing randomized trial blocks, CPU core pinning, nanosecond latency sampling, and OS resource telemetry via \texttt{getrusage} and system-call tracing.

\subsection{Userspace Communication Layer}
All six userspace IPC implementations adhere to a unified invocation signature: \texttt{ipc\_<mode>\_run\_opts(payload\_size, count, opts, \&elapsed\_ns, \&errors)}. The configuration structure \texttt{ipc\_bench\_opts\_t} allows caller-specified untimed warm-up cycles, explicit CPU affinity pinning for producer and consumer processes, and optional collection of per-exchange round-trip time (RTT) samples into pre-allocated memory buffers. 
To eliminate startup race conditions during process rendezvous, all mechanisms enforce strict synchronization protocols prior to entering the timed execution loop: processes establish channels, configure signal handlers, and execute warm-up exchanges before starting monotonic timing.

\subsection{Adaptive Transport Selection}
The adaptive framework addresses the trade-off between low-latency shared-memory synchronization for small messages and high-bandwidth kernel streaming for large messages. The system architecture supports two selection tiers:
\begin{itemize}
    \item \textbf{Static Size-Threshold Selector:} A deterministic, zero-overhead branch evaluating payload size $s$ against threshold $T$ (defaulting to 1,024~B, matching \texttt{SHM\_MAX\_PAYLOAD}):
    \begin{equation}
    \mathcal{T}_{\text{static}}(s) = 
    \begin{cases} 
    \text{SHM}, & \text{if } s \le \min(T, S_{\max}) \\ 
    \text{PIPE}, & \text{otherwise} 
    \end{cases}
    \end{equation}
    where $S_{\max} = 1024$~B is the maximum shared memory payload capacity.
    \item \textbf{Dynamic EWMA Telemetry Controller:} A runtime feedback controller that maintains rolling latency estimates for active transports, employing an 8\% hysteresis band to suppress transport thrashing under transient scheduling noise (detailed in Section~\ref{sec:adaptive_kernel}).
\end{itemize}

\subsection{Kernel Mailbox and Benchmarking Framework}
The kernel device mailbox (\texttt{/dev/aipc}) is an out-of-tree character device driver providing a kernel-buffered alternative to pipes and sockets. The driver maintains a dynamically allocated 64~KB circular ring buffer in kernel space. Reader and writer threads synchronize using Linux wait queues and a sleepable mutex. Crucially, the architecture incorporates an atomic open-count tracking mechanism and a teardown barrier in \texttt{aipc\_exit()} that mathematically eliminates module-unload race conditions.
To ensure statistical validity, the unified benchmarking suite (\texttt{bench/unified\_bench.c}) implements three controls: trial-block randomization via Fisher-Yates shuffling, dedicated CPU core pinning via \texttt{sched\_setaffinity()}, and high-resolution timing measured exclusively via \texttt{clock\_gettime(CLOCK\_MONOTONIC)}.
